\documentclass[11pt,oneside,openright,numbers=noenddot]{scrbook}
\usepackage{amsmath,amssymb,amsthm,placeins}
\usepackage[usenames,svgnames,dvipsnames]{xcolor}
\usepackage{hyperref}
\usepackage[nameinlink]{cleveref}
\usepackage[shortlabels]{enumitem}


\usepackage{graphicx,tikz}
\usepackage{caption}
\usepackage{tikz-cd}
\usepackage{pgfplots}
\pgfplotsset{compat=1.5}

\definecolor{azuree}{HTML}{007FFF}
\definecolor{Cobalt}{HTML}{0047AB}



\newcommand{\cbrt}[1]{\sqrt[3]{#1}}
\newcommand{\floor}[1]{\left\lfloor #1 \right\rfloor}
\newcommand{\ceil}[1]{\left\lceil #1 \right\rceil}
\newcommand{\mailto}[1]{\href{mailto:#1}{\texttt{#1}}}
\newcommand{\ol}{\overline}
\newcommand{\ul}{\underline}
\newcommand{\wt}{\widetilde}
\newcommand{\wh}{\widehat}
\newcommand{\eps}{\varepsilon}
\newcommand{\half}{\frac{1}{2}}



\DeclareMathOperator{\cis}{cis}
\DeclareMathOperator{\Int}{Int}
\DeclareMathOperator{\Bd}{Bd}
\DeclareMathOperator{\Lk}{Lk}
\DeclareMathOperator{\St}{St}
\DeclareMathOperator{\Sd}{Sd}
\DeclareMathOperator{\sd}{sd}
\DeclareMathOperator{\Ext}{Ext}
\DeclareMathOperator{\id}{id}
\DeclareMathOperator{\ran}{range}
\DeclareMathOperator{\diam}{diam}
\DeclareMathOperator*{\lcm}{lcm}
\DeclareMathOperator*{\argmin}{arg min}
\DeclareMathOperator*{\argmax}{arg max}



\newcommand{\CC}{\mathbb C}
\newcommand{\FF}{\mathbb F}
\newcommand{\NN}{\mathbb N}
\newcommand{\QQ}{\mathbb Q}
\newcommand{\RR}{\mathbb R}
\newcommand{\EE}{\mathbb E}
\newcommand{\ZZ}{\mathbb Z}
\newcommand{\PP}{\mathbb P}


\newcommand{\abs}[1]{\left\lvert #1 \right\rvert}
\newcommand{\norm}[1]{\left\lVert #1 \right\rVert}


\renewcommand{\qedsymbol}{$\blacksquare$}



\newenvironment{subproof}[1][Proof]{%
\begin{proof}[#1] \renewcommand{\qedsymbol}{$\square$}}%
{\end{proof}}



\newenvironment{soln}{\begin{proof}[Solution]}{\end{proof}}

\newcommand{\liff}{\leftrightarrow}
\newcommand{\lthen}{\rightarrow}
\newcommand{\opname}{\operatorname}
\newcommand{\surjto}{\twoheadrightarrow}
\newcommand{\injto}{\hookrightarrow}
\newcommand{\del}{\partial}

\renewcommand{\Re}{\opname{Re}}
\renewcommand{\Im}{\opname{Im}}

\DeclareMathOperator{\im}{im} % Image
\DeclareMathOperator{\Img}{Im} % Image
\DeclareMathOperator{\coker}{coker} % Cokernel
\DeclareMathOperator{\Coker}{Coker} % Cokernel
\DeclareMathOperator{\Ker}{Ker} % Kernel
\DeclareMathOperator{\rank}{rank}
\DeclareMathOperator{\dist}{dist}
\DeclareMathOperator{\Cok}{Cok} % Cokernel
\DeclareMathOperator{\cok}{cok} % Cokernel



\hypersetup{
    colorlinks,
    linkcolor={blue},
    citecolor={blue!50!black},
    urlcolor={blue!80!black}
}



% THESE ARE FOR HEADERS AND FOOTERS 


\usepackage[headsepline]{scrlayer-scrpage}
\renewcommand{\headfont}{}
\addtolength{\textheight}{3.14cm}
\setlength{\footskip}{0.5in}
\setlength{\headsep}{10pt}
%
\automark[chapter]{chapter}



\lohead{\footnotesize \leftmark}
\chead{}
\rofoot{}
\refoot{}
\lefoot{}
\lofoot{}


% 	THESE LINES INDICATE HOW THE CHAPTERS AND SECTIONS AND SUBSECTIONS LOOK LIKE
\renewcommand*{\sectionformat}{\color{blue}\S\thesection\autodot\enskip}
\renewcommand*{\subsectionformat}{\color{blue}\S\thesubsection\autodot\enskip}
\renewcommand{\thesubsection}{\thesection.\roman{subsection}}

\addtokomafont{chapterprefix}{\raggedleft}
\RedeclareSectionCommand[beforeskip=0.5em]{chapter}
\renewcommand*{\chapterformat}{%
\mbox{\scalebox{1.5}{\chapappifchapterprefix{\nobreakspace}}%
\scalebox{2.718}{\color{blue}\thechapter\autodot}\enskip}}



\usepackage{thmtools}
\usepackage[framemethod=TikZ]{mdframed}

\theoremstyle{definition}
\mdfdefinestyle{mdbluebox}{%
	roundcorner = 10pt,
	linewidth=1pt,
	skipabove=12pt,
	innerbottommargin=9pt,
	skipbelow=2pt,
	nobreak=true,
	linecolor=blue,
	backgroundcolor=TealBlue!5,
}
\declaretheoremstyle[
	headfont=\sffamily\bfseries\color{MidnightBlue},
	mdframed={style=mdbluebox},
	headpunct={\\[3pt]},
	postheadspace={0pt}
]{thmbluebox}

\mdfdefinestyle{mdredbox}{%
	linewidth=0.5pt,
	skipabove=7pt,
	frametitleaboveskip=5pt,
	frametitlebelowskip=0pt,
	skipbelow=2pt,
	frametitlefont=\bfseries,
	innertopmargin=4pt,
	innerbottommargin=8pt,
	% nobreak=true,
	linecolor=RawSienna,
	backgroundcolor=Salmon!5,
}
\declaretheoremstyle[
	headfont=\bfseries\color{RawSienna},
	mdframed={style=mdredbox},
	headpunct={\\[3pt]},
	postheadspace={0pt},
]{thmredbox}

	\mdfdefinestyle{mdneelbox}{%
			skipabove=8pt,
			linewidth=2pt,
			rightline=false,
			leftline=true,
			topline=false,
			bottomline=false,
			linecolor=RoyalBlue,
			backgroundcolor=RoyalBlue!5,
		}
		\declaretheoremstyle[
			headfont=\bfseries\sffamily\color{RoyalBlue!70!black},
			bodyfont=\normalfont,
			spaceabove=2pt,
			spacebelow=1pt,
			mdframed={style=mdneelbox},
			headpunct={\\[3pt]},
		]{thmneelbox}
	%%%%%%%%
		\mdfdefinestyle{mdazurebox}{%
			skipabove=8pt,
			linewidth=2pt,
			rightline=false,
			leftline=true,
			topline=false,
			bottomline=false,
			linecolor=azuree,
			backgroundcolor=azuree!5,
		}
		\declaretheoremstyle[
			headfont=\bfseries\sffamily\color{azuree!70!black},
			bodyfont=\normalfont,
			spaceabove=2pt,
			spacebelow=1pt,
			mdframed={style=mdazurebox},
			headpunct={\\[3pt]},
		]{thmazurebox}
		%%%%%%
		\mdfdefinestyle{mdcobaltbox}{%
			skipabove=8pt,
			linewidth=2pt,
			rightline=false,
			leftline=true,
			topline=false,
			bottomline=false,
			linecolor=Cobalt,
			backgroundcolor=Cobalt!5,
		}
		\declaretheoremstyle[
			headfont=\bfseries\sffamily\color{Cobalt!70!black},
			bodyfont=\normalfont,
			spaceabove=2pt,
			spacebelow=1pt,
			mdframed={style=mdcobaltbox},
			headpunct={\\[3pt]},
		]{thmcobaltbox}


\declaretheorem[%
style=thmcobaltbox,name=Theorem,numberwithin=chapter]{theorem}
\declaretheorem[style=thmazurebox,name=Lemma,sibling=theorem]{lemma}
\declaretheorem[style=thmneelbox,name=Proposition,sibling=theorem]{proposition}
\declaretheorem[style=thmbluebox,name=Corollary,sibling=theorem]{corollary}
\declaretheorem[style=thmredbox,name=Example,numberwithin=chapter]{example}

\mdfdefinestyle{mdgreenbox}{%
	skipabove=8pt,
	linewidth=2pt,
	rightline=false,
	leftline=true,
	topline=false,
	bottomline=false,
	linecolor=ForestGreen,
	backgroundcolor=ForestGreen!5,
}
\declaretheoremstyle[
	headfont=\bfseries\sffamily\color{ForestGreen!70!black},
	bodyfont=\normalfont,
	spaceabove=2pt,
	spacebelow=1pt,
	mdframed={style=mdgreenbox},
	headpunct={ --- },
]{thmgreenbox}

\mdfdefinestyle{mdkalabox}{%
	skipabove=8pt,
	linewidth=2pt,
	rightline=false,
	leftline=true,
	topline=false,
	bottomline=false,
	linecolor=black,
	backgroundcolor=RedViolet!5!gray!5,
}
\declaretheoremstyle[
	headfont=\bfseries\sffamily\color{black},
	bodyfont=\normalfont,
	spaceabove=1pt,
	spacebelow=1pt,
	mdframed={style=mdkalabox},
	headpunct={.},
]{thmkalabox}
\declaretheorem[style=thmkalabox,name=Definition,numberwithin=chapter]{defn}


\newtheorem*{abuse}{Abuse of Notation}
\newtheorem*{ques}{Question}
\newtheorem*{remark}{Remark}
\newtheorem{exercise}{Exercise}[chapter]
\newtheorem{problem}{Problem}[chapter]
\newtheorem{ex}{Example}[chapter]
\newtheorem{axiom}{Axiom}






\title{%
\includegraphics[scale=.75]{MAT215/tex/bracu_logo.png}~\\[2cm]
Complex Variables \& Laplace Transformations (MAT215)
}
\subtitle{ ~\\[1cm]
\LARGE{Lecture Notes}\\
\vspace{1cm}
\Large{Author: Shovan Sourav Datta Pranta (SVP)}\\
\Large{Editor: }\\
\Large{Template: Atonu Roy Chowdhury}}
\date{\today}
\usepackage{geometry}
\geometry{
a4paper,
total={170mm,257mm},
left=20mm,
top=20mm,
}

\begin{document}
\frontmatter

\begin{center}
    \Large \textbf{Formula Sheet} \\
    \Large MAT215: Complex Variables \& Laplace Transformations \\
    \Large BRAC University \\
\end{center}
\hrule 
\begin{enumerate}
    \item \textbf{Laplace Transformation of the function $f(t)$} : $$F(s)=\mathcal{L}\{f(t)\}=\int_0^\infty f(t)e^{-st} dt$$ 
    \item \textbf{Laplace Transformation table:}
    \begin{table}[ht!]
        \centering
        \renewcommand{\arraystretch}{2.25}
        \begin{tabular}{|c|c|}
            \hline
             $f(t)$ & $\mathcal{L}\{f(t)\}=F(s)$  \\
             \hline
            \hline
             $1$ & $\displaystyle \frac{1}{s}$ \hspace{0.5cm} $Re(s)>0$\\
             \hline 
             $t$ & $\displaystyle \frac{1}{s^2}$ \hspace{0.5cm} $Re(s)>0$\\
             \hline
             $t^n$; \hspace{0.1cm} $n=0,1,2,...$ & $\displaystyle \frac{n!}{s^{n+1}}$ \hspace{0.5cm} $Re(s)>0$\\
             \hline
             $e^{at}$ & $\displaystyle \frac{1}{s-a}$ \hspace{0.5cm} $Re(s)>a$\\
             \hline
             $\sin (at)$ & $\displaystyle \frac{a}{s^2+a^2}$ \hspace{0.5cm} $Re(s)>|a|$\\
             \hline
             $\cos (at)$ & $\displaystyle \frac{s}{s^2+a^2}$ \hspace{0.5cm} $Re(s)>|a|$\\
             \hline
             $\sinh (at)$ & $\displaystyle \frac{a}{s^2-a^2}$ \hspace{0.5cm} $Re(s)>|a|$\\
             \hline
             $\cosh (at)$ & $\displaystyle \frac{s}{s^2-a^2}$ \hspace{0.5cm} $Re(s)>|a|$\\
             \hline
        \end{tabular}
    \end{table}
    \FloatBarrier
    Here $s$ is a complex variable and $Re(s)$ indicates the real part of $s$. 
    \item \textbf{Laplace Transformations of derivatives:} 
    $$\mathcal{L} \{y^n(t)\}=s^n\mathcal{L} \{y(t)\}-s^{n-1}y(0)-s^{n-2}y'(0)- ......... -sy^{n-2}(0)-y^{n-1}(0)$$
    \item \textbf{First Translation Theorem:}
    If $F(s)=\mathcal{L} \{f(t)\}$, then
    \begin{itemize}
        \item $\displaystyle \mathcal{L} \{e^{at}f(t)\}=F(s-a)$
        \item $\displaystyle \mathcal{L}^{-1}\{F(s-a)\} =e^{at}f(t)$ \quad  or, \quad $\displaystyle \mathcal{L}^{-1}\{F(s)\} =e^{at}\mathcal{L}^{-1}\{F(s+a)\}$
    \end{itemize}
    \item \textbf{Unit Step Function:} $$u(t-a)=\begin{cases}
        0 & t \leq a\\
        1 & t>a
    \end{cases}$$
    \item \textbf{Second Translation Theorem:} If $F(s)=\mathcal{L} \{f(t)\}$ and $a>0$, then
    \begin{itemize}
        \item If $\displaystyle \mathcal{L} \{f(t-a)u(t-a)\}=e^{-as}F(s)$ 
        \quad or, \quad $\displaystyle \mathcal{L} \{f(t)u(t-a)\}=e^{-as}\mathcal{L} \{f(t+a)\}$
        \item $\displaystyle \mathcal{L}^{-1} \{e^{-as}F(s)\}=f(t-a)u(t-a)$
    \end{itemize}
    \item If $F(s)=\mathcal{L} \{f(t)\}$, then $\displaystyle \mathcal{L} \{t^n f(t)\}=(-1)^n \frac{d^n}{ds^n}\{F(s)\}$,\quad \quad for $n=1,2,3,...$
\end{enumerate}

\noindent \textbf{\Large Line Integral}
\begin{enumerate}
    \item \textbf{Complex line integral:}
    \begin{itemize}
        \item $ \displaystyle\int_a^b f(z) \, dz = \oint_C f(z) \, dz = \int_{t_1}^{t_2} f(z(t)) z'(t) \, dt \quad t_1 \leq t \leq t_2$
    \item  $\displaystyle \int_{t_1}^{t_2} f(z(t)) z'(t) \, dt = \int_a^b \sqrt{[x'(t)]^2 + [y'(t)]^2} \, dt$
    
    \end{itemize}
    
    \item \textbf{Path Independence:}\\
    If \( f(z) \) has an antiderivative in an open region \(A\), then the path integral \(\displaystyle \oint_C f(z) \, dz \) is path independent for all paths in \(A\).
    
    The following two things are equivalent:
    \begin{itemize}
        \item The integral \(\displaystyle  \oint_C f(z) \, dz \) is path independent.
        \item The integral \(\displaystyle  \oint_C f(z) \, dz \) around any closed path is 0.
    \end{itemize}
\end{enumerate}

\backmatter


\end{document}
